\let\negmedspace\undefined
\let\negthickspace\undefined
\documentclass[journal]{IEEEtran}
\usepackage[a5paper, margin=10mm, onecolumn]{geometry}
%\usepackage{lmodern} % Ensure lmodern is loaded for pdflatex
\usepackage{tfrupee} % Include tfrupee package

\setlength{\headheight}{1cm} % Set the height of the header box
\setlength{\headsep}{0mm}     % Set the distance between the header box and the top of the text

\usepackage{gvv-book}
\usepackage{gvv}
\usepackage{cite}
\usepackage{amsmath,amssymb,amsfonts,amsthm}
\usepackage{algorithmic}
\usepackage{graphicx}
\usepackage{textcomp}
\usepackage{xcolor}
\usepackage{txfonts}
\usepackage{listings}
\usepackage{enumitem}
\usepackage{mathtools}
\usepackage{gensymb}
\usepackage{comment}
\usepackage[breaklinks=true]{hyperref}
\usepackage{tkz-euclide} 
\usepackage{listings}
% \usepackage{gvv}                                        
\def\inputGnumericTable{}                                 
\usepackage[latin1]{inputenc}                                
\usepackage{color}                                            
\usepackage{array}                                            
\usepackage{longtable}                                       
\usepackage{calc}                                             
\usepackage{multirow}                                         
\usepackage{hhline}                                           
\usepackage{ifthen}                                           
\usepackage{lscape}
\begin{document}

\bibliographystyle{IEEEtran}
\vspace{3cm}

\title{12.47}
\author{EE25BTECH11033 - Kavin}
% \maketitle
% \newpage
% \bigskip
{\let\newpage\relax\maketitle}

\renewcommand{\thefigure}{\theenumi}
\renewcommand{\thetable}{\theenumi}
\setlength{\intextsep}{10pt} % Space between text and floats
\textbf{Question}:\\
The following simultaneous equations
\begin{align*}
x + y + z &= 3 \\
x + 2y + 3z &= 4 \\
x + 4y + kz &= 6
\end{align*}
will NOT have a unique solution for $k$ equal to 
\begin{enumerate}
\begin{multicols}{4}
\item $0$
\item $5$
\item $6$
\item $7$
\end{multicols}
\end{enumerate}\\    
\bigskip

\textbf{Solution}:\\
The equations can be written as,
\begin{align}
    \myvec{1 & 1 & 1 \\ 1 & 2 & 3 \\ 1 & 4 & k }\vec{x}=\myvec{3\\4\\6}
\end{align}
where
\begin{align*}
    \vec{x} = \myvec{x\\y\\z}
\end{align*}
The augmented matrix is given by,
\begin{align}
    \augvec{3}{1}{1 & 1 & 1 & 3 \\
1 & 2 & 3 & 4 \\
1 & 4 & k & 6}
\end{align}
\begin{align}
    R_2 \rightarrow R_2 - R_1\text{ and }R_3 \rightarrow R_3 - R_1\implies \augvec{3}{1}{1 & 1 & 1 & 3 \\
0 & 1 & 2 & 1 \\
0 & 3 & k-1 & 3}
\end{align}
\begin{align}
    R_3 \rightarrow R_3 - 3R_2\implies\augvec{3}{1}{1 & 1 & 1 & 3 \\
0 & 1 & 2 & 1 \\
0 & 0 & k-7 & 0}
\end{align}
\begin{align}
    \myvec{1 & 1 & 1\\0 & 1 & 2\\0 & 0 & k-7}\vec{x}=\myvec{3\\1\\0}
\end{align}
For the system to \textbf{NOT} have a unique solution, the coefficient of $z$ in the final row must be zero.
\begin{align}
    k-7=0   
\end{align}
\begin{align}
    \implies k=7
\end{align}

\begin{figure}[H]
\begin{center}
\includegraphics[width=0.8\columnwidth]{figs/fig.png}
\end{center}
\label{fig:Fig1}
\end{figure}
\end{document}


